\section{Cell paths}\label{sec:geometry}


We now construct paths between distant cells. Their purpose is to keep the
second moment of an averaged endpoint change bounded as the volume grows.
The first step controls encounters of two random skeletons; the second
step adapts each skeleton to the actual unusable cells.

\subsection{Random skeletons and their encounter moments}\label{sec:skeletons}

The use of unpredictable paths and intersection moments follows
Benjamini--Pemantle--Peres~\cite[Theorem~1.3 and Lemma~3.1]{BenjaminiPemantlePeres1998}.
The independent block velocities below are a continuous dyadic version
of the aligned block-drift construction of H\"aggstr\"om and
Mossel~\cite[Section~3, Proposition~3.1]{HaggstromMossel1998}.
Finite-endpoint path averaging appears in the synchronization argument
of Abbe, Massouli\'e, Montanari, Sly, and
Srivastava~\cite[Section~6, Lemma~6.1]{AbbeMassoulieMontanariSlySrivastava2018}
and in the symmetry-breaking argument of Garban and
Spencer~\cite[Section~2, Lemma~2.5]{GarbanSpencer2023}.
We give an explicit dyadic construction and prove the thickened and
soft-distance estimates needed here.

The construction in this subsection is purely geometric. Let
$\mathcal G_n=(\Z/n\Z)^3$, and fix an ordered pair $a,b$ with lifted
displacement $(w,d_2,d_3)$ satisfying
\[
 \lceil n/10\rceil\le w\le\lfloor n/8\rfloor,
 \qquad |d_2|,|d_3|\le w.
\]
Call such a pair admissible. Admissible pairs form a fraction at least
an absolute $\upsilon>0$ of all ordered pairs for sufficiently large $n$.
All coordinates in the construction are first taken relative to a fixed
lift of $a$, and then projected to $\mathcal G_n$.

The first coordinate will advance from $0$ to $w$. In the other two
coordinates we add a random displacement to the straight interpolation
between the endpoints. We read that displacement forward up to the
midpoint and backward afterward, so it vanishes at both endpoints.
This fixes the endpoint without conditioning the random increments;
the two halves deliberately use the same displacement variables.

For $q=2,3$, let
$(U^{(q)}_{\ell,j})_{\ell,j\ge0}$ be mutually independent random
variables, uniform on $[-1,1]$. Fix
$c_{\rm vel}=(1-2^{-1/4})/2$, and set
\[
 v_i^{(q)}=c_{\rm vel}\sum_{\ell\ge0}2^{-\ell/4}
          U^{(q)}_{\ell,\lfloor(i-1)/2^\ell\rfloor},
 \qquad Z_t^{(q)}=\sum_{i=1}^t v_i^{(q)},\qquad Z_0^{(q)}=0.
\]
The series converges absolutely, and $|v_i^{(q)}|\le1/2$.
At integer times $t=0,\ldots,w$ use the base vertices
\[
 \left(t,\left\lfloor td_2/w+Z^{(2)}_{\min(t,w-t)}\right\rfloor,
          \left\lfloor td_3/w+Z^{(3)}_{\min(t,w-t)}\right\rfloor\right).
\]
Starting from the base vertex at time $t$, first advance one step in
the first coordinate, then take the signed second-coordinate steps,
and finally the signed third-coordinate steps needed to reach the
base vertex at time $t+1$. Omit zero-length moves. Begin the vertex
list with $a$ and append each newly reached vertex once, without
repeating a shared endpoint of successive pieces. Each transverse
coordinate changes by at most two; thus the resulting
nearest-neighbor skeleton $\gamma=(z_0,\ldots,z_l)$ has at most five
steps per time interval and at most five listed occurrences at each
first-coordinate level $t\ge1$, with one occurrence at level zero.
The law depends only on the torus and the ordered endpoints.
Independent skeletons use independent arrays of the $U$'s.
The dyadic weights are chosen so that a previously unused block of
$k$ increments supplies a transverse displacement of order $k^{3/4}$.
In two transverse coordinates this gives the summable conditional
density bound $Ck^{-3/2}$ proved below.

For vertices on the torus write
\[
 d_\infty(x,y)=\min_{k\in\Z^3}\|\widetilde x-\widetilde y+nk\|_\infty,
 \qquad d_\infty(x,\gamma')=\min_{z'\in\gamma'}d_\infty(x,z'),
\]
where the tildes denote arbitrary lifts. For two independent skeletons
$\gamma,\gamma'$ with the same ordered endpoints and a fixed radius
$R_*\ge1$, define
\[
 J_0=\sum_{z\in\gamma}\mathbf1_{\{d_\infty(z,\gamma')\le R_*\}}
       +\sum_{z'\in\gamma'}\mathbf1_{\{d_\infty(z',\gamma)\le R_*\}},
 \qquad
 S_B=\sum_{z\in\gamma}e^{-d_\infty(z,\gamma')/B}\quad(B>0).
\]
Every sum uses the listed vertex occurrences, including the endpoints.
An occurrence contributes once to $J_0$ when at least one occurrence
of the other skeleton lies within the stated distance.

\begin{lemma}[Skeleton encounters]\label{lem:skeleton}
For the independent skeletons just constructed, there are absolute
constants $c_{\mathrm{hit}}>0$ and $C<\infty$ such that, for every
$R_*\ge1$,
\[
             \mathbb E e^{\alpha J_0}\le C,\qquad
             \alpha=\frac{c_{\rm hit}}{(1+R_*)^2}.
             \tag{24}\label{eq:24}
\]
For every fixed $B>0$ there are $\xi(B)>0$ and $C(B)<\infty$ such that
\[
             \mathbb E e^{\xi(B)S_B}\le C(B).
             \tag{25}\label{eq:25}
\]
Both estimates are uniform in $n$ and in every admissible ordered pair.
\end{lemma}
\begin{proof}
\emph{Conditional density bound.}
In this proof write
\(Z_t=(Z_t^{(2)},Z_t^{(3)})\) for the pair of transverse sums,
and write \(Z'_t\) for the corresponding pair of the other skeleton.
Put \(h=\lfloor w/2\rfloor\). Condition on all the variables
generating \(\gamma'\). For the first skeleton define the dyadic blocks
\[
 I_{\ell,j}=\{j2^\ell+1,\ldots,(j+1)2^\ell\},
\]
and let \(\mathcal F_t\) be generated by the entire primed array and
the variables \(U^{(q)}_{\ell,j}\) for which
\(I_{\ell,j}\cap\{1,\ldots,t\}\ne\varnothing\), with \(q=2,3\).
Thus it reveals exactly the variables of the first array used through time
\(t\), including the large blocks that continue past that time.
For \(t+k\le h\), \(k\ge1\), the future index interval contains
a full aligned dyadic block of length at least \(k/4\): take a
power of two between \(k/4\) and \(k/2\) and skip less than one
such block to align, or use length one for \(k=1\). Its variable
has not been used and contributes a uniform term of width at least
\(c k^{3/4}\) to the sum, even after conditioning on the other
variables. Indeed, a block at level \(\ell\) contributes
\(c_{\rm vel}2^{3\ell/4}\) times its variable. Using one such
variable in each independent coordinate gives a conditional
two-coordinate density of \(Z_{t+k}\), given \(\mathcal F_t\),
bounded by \(Ck^{-3/2}\).

\emph{Same-time encounter tests.}
Write \(D_t=|Z_t-Z'_t|_\infty\), for \(0\le t\le h\).
Denote the two base vertices at time \(t\) by \(b_t,b'_t\).
An occurrence \(z\) at first-coordinate level \(t\) is within
sup-norm distance two of \(b_t\). If \(z'\) is an occurrence of
the other skeleton at level \(t'\), then
\( |t-t'|\le d_\infty(z,z')\): first-coordinate differences
are at most \(w\le n/8\), so they do not wrap. Each base step
has sup-norm size at most two. The triangle inequality therefore gives
\[
 d_\infty(b_t,b'_t)\le 3d_\infty(z,z')+4.
\]
The same-time lifted transverse differences also do not wrap.
The linear drifts are common, each random deviation has absolute
value at most \(w/4\), and rounding changes their difference by
at most one. Thus these differences have magnitude at most
\(w/2+1<n/2\). With \(u=\min(t,w-t)\), it follows that
\[
 D_u\le 3d_\infty(z,z')+5.
\]

Consequently
\[
 J_0\le C\sum_{t=0}^{h}
             \mathbf1_{\{D_t\le C(R_*+1)\}}.
\]
The reflected times \(t\) and \(w-t\) have the same deviation
\(Z_t\); their different linear drifts are common to both skeletons
and cancel in the difference. Thus the second half contributes only
another fixed multiplicity to this bound.

For the distance-weighted sum, take \(z'\) to be a closest
occurrence of \(\gamma'\). The same inequality gives
\[
 e^{-d_\infty(z,\gamma')/B}\le
 e^{5/(3B)}e^{-D_u/(3B)},\qquad
 S_B\le C(B)\sum_{t=0}^{h}e^{-D_t/(CB)}.
\]
These are the two deterministic reductions used in the moment estimate.

\emph{Exponential moment.}
Let \(X_t\in[0,1]\) denote either of these first-half tests.
They are adapted to \(\mathcal F_t\), since all variables of
\(\gamma'\) were conditioned on at the outset. Integrating the
conditional density bound against the indicator, respectively
the decaying exponential, and summing over the positive time
increments gives the uniform estimate
\[
 \mathbb E\left[\sum_{u=t+1}^{h}X_u\,\middle|\,\mathcal F_t\right]
 \le b_0,
 \qquad
 b_0=\begin{cases}
 C(R_*+1)^2,&X_u=\mathbf1_{\{|Z_u-Z'_u|_\infty\le C(R_*+1)\}},\\
 C(B),&X_u=e^{-|Z_u-Z'_u|_\infty/(CB)}.
 \end{cases}
\]
For every \(j\ge1\), the expected sum of products on strictly
ordered \(j\)-tuples is therefore at most \((b_0+1)^j\).
To see this, condition the last-index sum on the revealment at
each preceding index in the nested sum, and include \(X_0\le1\)
when bounding the one-term sum. Now expand
\[
 e^{\theta\sum_{t=0}^{h}X_t}
 =\prod_{t=0}^{h}\bigl(1+(e^{\theta X_t}-1)\bigr).
\]
The inequality \(e^{\theta x}-1\le(e^\theta-1)x\) for
\(0\le x\le1\) bounds its expectation by the geometric series
\(\sum_{j\ge0}[(e^\theta-1)(b_0+1)]^j\).
Taking \(\theta\) to be a sufficiently small constant times
\(1/(b_0+1)\), and accounting for the preceding sum-comparison
constants, proves (24) and (25).
\end{proof}

\subsection{Detours around unusable components}

We now apply the geometric construction to the coarse grid of
Section~\ref{sec:setup}. This is the point at which the physical cutoffs
enter: set
\[
 R_* = \lceil20(R+R_A+D+20)\rceil,\qquad
 \varepsilon=\alpha/20,
\]
decreasing the absolute constant in $\alpha$ if necessary.
Choose $D$ large enough for the earlier conditions with this choice.
This is possible because \eqref{eq:17} has Gaussian decay in $D$,
whereas $R_*$ grows only linearly. Choose $\zeta,C_0$ and then the
small-density threshold as in Proposition~\ref{prop:defective-cells}.
Increase $C_0$ also so that, for two uniform anchors, the probability
of any end-score failure at threshold $C_0$ is less than $\upsilon/8$;
the spatial average after \eqref{eq:16} gives this last requirement.

The skeleton law does not see the unusable cells. We next adapt each
sampled skeleton to those cells by a deterministic detour procedure.
This step uses nearest-neighbor components of bad cells but allows
sup-norm-one steps on the resulting path. Here and below, bad means
unusable in the sense of Section~\ref{sec:slab}. Reject the configuration
if any nearest-neighbor component of bad cells has size \(\ge n/4\).

A smaller component \(T\) lifts consistently to a finite connected
lattice set, up to translation. Indeed, simple cycles in it are too
short to wrap, and closed walks decompose into simple cycles and
backtracks. Its bounding box has coordinate diameters less than
\(n/4\), since two of its vertices can be joined within \(T\) by
fewer than \(n/4\) steps. Define its \emph{fill} by adjoining the
finite nearest-neighbor components of its lattice complement, and
project the result back to the torus. The fill lies in the same
bounding box: a vertex outside that box connects straight outwards
to infinity without meeting \(T\). Reject if either endpoint lies
in any projected fill.

The detour procedure scans an accepted skeleton. On entering an
outermost fill, meaning one not contained in another fill, it replaces
the portion through its last visit there by an exterior-boundary path.
The next proof establishes the boundary
connectivity that permits this replacement with sup-norm-one steps.
Fix all tie-breaking rules for boundary paths in advance, and use
chronological loop erasure for the resulting walk.

\begin{lemma}[Deterministic detours]\label{lem:detour}
Apply the rejection rules just specified to a nearest-neighbor skeleton
from $a$ to $b$ and a set of bad cells in $(\Z/n\Z)^3$. If the skeleton
is accepted, the procedure below produces a simple path with sup-norm
steps of length one from $a$ to $b$, using only good cells. Every vertex
is a skeleton vertex or lies on the exterior vertex boundary of a bad
component hit by the skeleton. A component $T$ is detoured at most once
and contributes at most $6|T|$ detour vertices. Removing any whole bad
components not hit by an accepted skeleton preserves acceptance and its
output path.
\end{lemma}
\begin{proof}
\emph{Fills and projection.}
We first establish the boundary facts needed by the procedure.
In the lattice the fill is connected, since each hole neighbors
the obstacle. Its exterior is connected, since a finite set has
just one infinite complementary component. An edge from the fill
to its exterior starts in the obstacle itself.

For disjoint non-neighboring connected obstacles, their fills
are nested or disjoint and non-neighboring. Indeed, each obstacle
is in a single complementary component of the other. If it is in
a hole, its fill stays there: outside that hole one connects to
infinity avoiding the obstacle in it, either in the exterior of
the enclosing obstacle or by using that connected enclosing
obstacle, also to escape from any of its other holes. If the two
obstacles are mutually in the exteriors, each connected fill
avoids the other obstacle and thus stays in the other's exterior.
Adjacency would imply adjacency of the obstacles themselves by
the crossing-edge property.

These facts survive projection to the torus. An intersection or
adjacency can be lifted with corresponding translates of the two
obstacles and their fills. The small box size avoids self-overlap
and preserves cardinality, so proper nesting stays proper. We
therefore use the outermost fills in the construction below.

\emph{Connected exterior boundary.}
The exterior vertex boundary consists of the exterior neighbors
of the obstacle. It has at most \(6|T|\) vertices. We claim that
it is connected in the star graph, whose steps have sup-norm one.

We give a self-contained cycle-space proof of this boundary
connectivity fact; see Tim\'ar~\cite[arXiv version~2,
Lemmas~1--2]{Timar2013}. Consider the edge cut from the
lattice fill to its exterior. Link two cut edges when they belong
to a common elementary square. All cut edges belong to one link
component. Otherwise choose one link component. Every square
contains an even number of its edges: it contains an even number
of edges of the full cut, and all those edges are linked. The
same parity holds for every lattice cycle, since squares generate
cycles modulo two. Explicitly, in a closed walk commute steps in
different axes, each swap changing the walk by an elementary
square modulo two, and then cancel backtracks. Counting the
selected edges modulo two along paths from a fixed vertex now
defines a vertex potential whose differing-neighbor edges are
exactly the selected edges. This potential is constant on the
connected fill and on its connected exterior. Its cut must
therefore be either empty or the full cut, contradicting the
choice of a proper nonempty link component. This proves the claim
about cut edges. Their linked exterior endpoints are star
neighbors or identical, proving boundary connectivity.

Project this connected boundary back to the torus. Any edge
leaving the projected fill lifts from within the fill to outside
and hence goes from the projected obstacle to this projected
boundary. The boundary vertices are good: the small box size
places them outside the projected fill, and if one were bad it
would belong to the adjacent bad component. Thus each outermost
fill has a connected star boundary of at most \(6|T|\) usable
vertices, and every nearest-neighbor crossing uses that boundary.
This is the boundary information needed by the algorithm.

\emph{Detour construction.}
Scan the nearest-neighbor skeleton starting outside the outermost
fills. If the next vertex enters one, replace the portion through
the \emph{last} skeleton vertex in that same fill by a
star-boundary path from the exterior predecessor to the exterior
successor. These are boundary vertices by nearest-neighbor
crossing, and the successor is outside all outermost fills by
non-adjacency. Choose a deterministic simple path on the projected
boundary, with fixed tie-breaking depending only on that component
and the two vertices. Projection permits this connection regardless
of the endpoint lifts used. Continue scanning from the successor.

Finally apply chronological loop erasure to obtain a simple star
path \(P\) from \(a\) to \(b\) in good cells. Each detoured
component contributes at most \(6|T|\) detour vertices, all
adjacent to it. The last-visit rule ensures that each component
is detoured at most once; loop erasure can only remove vertices.

\emph{Deletion invariance.}
Suppose the original configuration is accepted. Any fill that
intersects the skeleton must have its obstacle hit by the
skeleton: the start is exterior and the skeleton has
nearest-neighbor steps. Enclosing fills of such obstacles are
hit as well. The construction is therefore determined by the
actual bad components intersecting the skeleton. Deleting other
\emph{whole} bad components in this accepted configuration
preserves acceptance and leaves both the detours and the final
loop-erased path unchanged.
\end{proof}

Rejection by large components has closed probability tending to zero at large \(n\) by (22): count connected sets of size \(\lceil n/4\rceil\) by at most \(M C^{\lceil n/4\rceil}\) (walk a rooted spanning tree), and a larger component contains one. For either given endpoint, containment in a small fill has probability at most \(\sum_{j\ge1} C(j+1)^3 C^j q_{\rm cl}^j\), since its bad component of size \(j\) must be within distance \(j\). This can be made arbitrarily small by the density restriction.
